\documentclass[10pt,a4paper]{amsart}
\usepackage[margin=19mm]{geometry}
\usepackage{amsmath,amssymb,mathtools}
\usepackage[scr=rsfs,cal=euler]{mathalfa}
\usepackage{xfrac}
\usepackage[unicode,hidelinks,bookmarks=false]{hyperref}
\newcommand{\ie}{i.e.\ }
\newcommand{\cf}{cf.\ }
\newcommand{\resp}{respectively\ }
\newcommand{\Subsection}{\subsection}
\providecommand{\nobreakdash}{\nobreak\hbox{-}}
\setlength{\emergencystretch}{2em}
\multlinegap0pt
\newtheorem{prop}{Proposition}
\title{$\ \ $ Structure and realizability for rational maps}
\author{Zhiqiang Wei}
\date{Source: arXiv:2511.06784v2 (CC BY 4.0)}
\begin{document}
\maketitle

\noindent\emph{Source and licence.} $\ \ $ Structure and realizability for rational maps. Version arXiv:2511.06784v2 (CC BY 4.0), distributed by arXiv under the Creative Commons Attribution 4.0 International licence, http://creativecommons.org/licenses/by/4.0/, as recorded in the arXiv licence metadata for this exact version and shown on the abstract page https://arxiv.org/abs/2511.06784v2. Full licence notice: \texttt{licenses/arxiv-2511.06784-CC-BY-4.0.txt}.\par\smallskip

\noindent\textit{Editorial scope.} Counted unit: the article's prop Pro-2-4 (source0.tex lines 378--384). The article's complete original proof of it is delivered with it (source0.tex lines 386--441, one \texttt{proof} environment of 56 lines). No mathematics is written, repaired or re-derived: the statement and the proof are the article's own lines, in the article's own notation, and no retained line is substituted or re-flowed.  Retained uncounted as the source-local context the proof needs: lines 366--368.  Nothing was omitted from any retained span, so the package carries no omission record.  No retained line contains a citation, so the wrapper carries no bibliography and no bibliography entry.  No retained line contains a figure environment, an image-inclusion command or drawing code, so no image is delivered and the package declares no asset dependency.  Editorial formatting only, in addition: an AMS \texttt{amsart} preamble that restates the article's own declarations and macros used by the retained lines, namely the environments \texttt{prop} on separate counters, together with the standard packages those lines need.  The retained environments print in this document as Proposition 1 in order of appearance, whereas the article prints the same items with the numbering of its own full document.  The complete native source of the article as distributed in the arXiv e-print for this exact version is bundled unchanged as \texttt{sources/188759/source0.tex}.
\begin{prop}\label{Pro-2-3}
In a small neighborhood of any extremum of $\Phi$, every integral curve of $\vec{V}$ is a topological circle enclosing the extremum in its interior.
\end{prop}

\begin{prop}\label{Pro-2-4}
The set $\operatorname{Sing}(\vec{V})$ decomposes into two disjoint subsets, $\operatorname{Sing}(\vec{V}) = S_1 \cup S_2$, where:
\begin{enumerate}
\item $S_1 = \{ p \in \operatorname{Sing}(\vec{V}) \mid \Omega_p = \emptyset \}$, and if $p \in S_1$, then $p$ is an extremum of $\Phi$.
\item $S_2 = \{ p \in \operatorname{Sing}(\vec{V}) \mid \Omega_p \neq \emptyset \}$, and if $p \in S_2$, then $p$ is a saddle point of $\Phi$ and the angle of $\mathrm{d}s^2$ at $p$ is $\pi \cdot |\Omega_p|$.
\end{enumerate}
\end{prop}

\begin{proof}
Case (1) follows from Proposition \ref{Pro-2-3}. We now prove Case (2). Without loss of generality, near a saddle point $p$ of $\Phi$, choosing a new complex coordinate $z$, we can assume $\omega = z^n \mathrm{d}z$ with $n \in \mathbb{Z}^{+}$, where $z$ is a local complex coordinate such that no other singularities of $\vec{V}$ lie in the neighborhood. Then $Y = \frac{\sqrt{-1}}{z^n} \frac{\partial}{\partial z}$. Writing $z = r e^{\sqrt{-1}\theta}$, we compute:
\begin{align*}
Y &= \frac{\sqrt{-1}}{r^n e^{\sqrt{-1}n\theta}} \left( \frac{\partial r}{\partial z} \frac{\partial}{\partial r} + \frac{\partial \theta}{\partial z} \frac{\partial}{\partial \theta} \right)
= \frac{\sqrt{-1}}{2 r^n e^{\sqrt{-1}(n+1)\theta}} \left( \frac{\partial}{\partial r} - \frac{\sqrt{-1}}{r} \frac{\partial}{\partial \theta} \right) \\
&= \frac{1}{2 r^n} \left[ \sin((n+1)\theta) \frac{\partial}{\partial r} + \frac{\cos((n+1)\theta)}{r} \frac{\partial}{\partial \theta} + \sqrt{-1}(\cdots) \right].
\end{align*}
Hence, the real vector field $\vec{V}$ is given by:
\[
\vec{V} = \frac{\sin((n+1)\theta)}{2 r^n} \frac{\partial}{\partial r} + \frac{\cos((n+1)\theta)}{2 r^{n+1}} \frac{\partial}{\partial \theta}.
\]

Suppose an integral curve of $\vec{V}$ near $p$ is parametrized as:
\[
\begin{cases}
r = r(t), \\
\theta = \theta(t),
\end{cases}
\quad t \in (-\varepsilon, \varepsilon) \setminus \{0\},\ \varepsilon > 0,
\]
with $\lim\limits_{t \to 0} r(t) = 0$. Then the system becomes:
\begin{equation}\label{P-2-E-1}
\begin{cases}
r' = \dfrac{\sin((n+1)\theta)}{2 r^n}, \\[10pt]
\theta' = \dfrac{\cos((n+1)\theta)}{2 r^{n+1}}.
\end{cases}
\end{equation}

If $\theta' \neq 0$, the solution of \eqref{P-2-E-1} satisfies:
\[
r^{-(n+1)} = \lambda |\cos((n+1)\theta)|, \quad \lambda > 0,
\]
i.e.,
\[
r = \frac{1}{\sqrt[n+1]{\lambda |\cos((n+1)\theta)|}} \geq \frac{1}{\sqrt[n+1]{\lambda}},
\]
so such a curve does not reach $p$.

If $\theta' = 0$, the solutions are:
\[
\begin{cases}
r^{n+1} = -\dfrac{n+1}{2} t, \\
\theta = \dfrac{1}{n+1} \left[ \dfrac{\pi}{2} + (2k+1)\pi \right], \quad k = 0, \ldots, n,
\end{cases}
\quad t < 0,
\]
and
\[
\begin{cases}
r^{n+1} = \dfrac{n+1}{2} t, \\
\theta = \dfrac{1}{n+1} \left( \dfrac{\pi}{2} + 2k\pi \right), \quad k = 0, \ldots, n,
\end{cases}
\quad t > 0.
\]
This completes the proof of Case (2).
\end{proof}

\end{document}
